\section{타당성 위협과 한계}
\label{sec:threats}

\subsection{구성 타당성}
CoC는 인간이 직접 작성한 안전 명세가 아니라 기록 궤적에 조건화된 VLM/LLM의 사후 주석이다. 따라서 문장이 장면의 원인이나 운전 의도를 충실하게 설명한다는 보장은 없다. \texttt{quality\_pass}, 파서와 의도 매핑은 감사 입력을 구성하는 기계 생성 신호일 뿐 인간 정답을 대체하지 않으며, 본 결과는 일반적인 CoC 충실도가 아니라 구성된 시간ㆍ출처 계약의 내부 정합성만 다룬다.

\subsection{측정 타당성}
응답은 액추에이터 명령이 아니라 자차 궤적에서 계산한 속도 변화 대리변수이다. 창 $W$, 속도 임계값 $\delta$, 방향 비율 $\rho$에 따라 판정이 달라질 수 있고, 27개 설정도 모든 유효 탐지기를 포괄하지 않는다. $v\leq0.5$~m/s STOP/YIELD 이미 충족 규칙은 초기 검토 후보 8건 중 3건을 재분류한 분석 대상 집합 보정 탐색 논리이다. 또한 $D=3$초는 실행 가능성 분석 설정이며 규범적 안전 기한이 아니다. 기한은 $r<D$로 해석하지만 저장 집계에는 정확히 $r=D$인 사례가 없고, 등호 시점 전환 우선순위는 후속 검사기 회귀 시험이 필요하다. \preserve/\reset 반복 정책과 기한 선택에 따른 판정 반전 때문에 계약 통과와 검토 후보는 측정ㆍ정책 조건을 함께 명시해야 한다.

\subsection{계약 타당성}
의도별 응답 술어, 후보 기한, 반복 사건 및 종료 시점 의미론은 연구자가 선택한 감사 계약이다. 이는 법규, 제어기 요구사항 또는 차량 동역학에서 독립적으로 도출한 규범 계약이 아니다. STOP/HOLD/RECOVER 다단계의 일부 전환도 합성 가정이므로 관찰자 표현력 시험에는 쓸 수 있지만 실제 행동 의무나 학습 정답으로 해석할 수 없다.

\subsection{내적 타당성}
일괄 결과는 파서, 품질 관문, 의도 매핑, 응답 탐지기와 결과 라우터의 구현에 의존한다. 282개 코호트 모델은 위반 불리언을 직접 설정하고 1,128개 기대 판정도 레이블에서 도출하므로 그 일치는 플래그/질의 직렬화 회귀만 지지한다. 관찰자 수준 시간 변이 근거는 \texttt{scene-0001} 모델에 한정하며 미관측 결함의 검출 성능은 측정하지 않는다. 특히 \texttt{orphan\_response}는 합성 비원자적 중복 제거 조건이며 자연 발생 오류의 관측치가 아니다.

\subsection{외적 타당성}
단일 장면 분석은 chunk~0의 94개 장면과 403개 주석 중 Qwen/Qwen3.5-35B-A3B의 자기평가를 통과하고 영어 정규식/복합 의도 규칙에서 방향이 명시된 134개 주석 사건별 후보 계약에 한정된다. 품질 통과 290개 중 모호ㆍ비방향ㆍ복합 의도 156개는 제외했으므로 추론은 이 부분집합에만 적용된다. 다른 nuScenes 분할, CoC 생성기, 언어, 횡방향 의도, 센서ㆍ차량 및 학습 정책으로 일반화하려면 별도 평가가 필요하다.

\subsection{검토 타당성}
다섯 후보의 전방 카메라 프레임 검토는 독립적인 인간 평가가 아니라 비인간 AI 시각 선별이었다. 네 프레임의 보조 문맥은 검토 경로를 정하는 데 사용할 수 있지만, 사건의 인과 관계, CoC 사실성 또는 물리적 위험에 대한 정답을 제공하지 않는다. 이후 연구에는 가려진 다중 평가자 검토와 합의 절차가 필요하다.

\subsection{출처 타당성}
로컬 CoC-Nusc 파일은 장면 간 로그ㆍ이전/다음 관계와 객체 추적 연속성을 제공하지 않으므로 기본적으로 독립 에피소드로 닫았다. 13체인 스모크 산출물은 공식 trainval 원본의 바이트 수준 출처를 재검증하지 못한 nuScenes 호환 메타데이터 미러에 의존한다. 더구나 느린 의도 우선 부분문자열 파서와 $W=1.0$초, $\delta=0.7$~m/s, 방향 비율 없는 탐지기는 공유 계약과 다르고 복합 의도를 오분류할 수 있다. 따라서 4/9 및 8/7/0은 해당 구현의 기술 출력일 뿐 공식 데이터, 객체ㆍ위험 연속성 또는 장기 계약 성능의 정량 근거가 아니다.

\subsection{정형화 한계}
현재 모델은 기록된 결정적 사건열과 자차 궤적을 재생하며 환경이나 제어 대안을 탐색하지 않는다. 집계 산출물이 13/13 일치를 기록하지만 실행 가능한 기준선과 항목별 출력이 없어 독립 재현성을 입증하지 않는다. 시간 관찰자는 모델군별 상태 생명주기, 교착과 반례 추적을 명시하지만 상대 객체 상태, 액추에이터, 차량 동역학 및 폐루프 상호작용을 포함하지 않으므로 물리 안전성, 제어 합성 또는 배포 적합성을 증명하지 않는다.
